\section{Proofs}
\label{app:proofs}

The theorem and corollary statements remain in the main text. Their proofs are collected
here so that the main narrative can emphasize the representation chain, geometric
controls, and numerical evidence.

\subsection{Spectral projection and finite hopping support}
\label{app:proof-projection-range}

\begin{proof}[Proof of Proposition~\ref{prop:projection-destroys-range}]
Set $z=e^{ik}$. If $d(k)$ were a finite Fourier polynomial, there would be a Laurent
polynomial $f(z)$ whose restriction to the unit circle equals $d(k)$. Squaring on the
unit circle and using equality on infinitely many points would give the Laurent-polynomial
identity
\begin{equation}
 f(z)^2=\Delta^2+2t^2+t^2(z+z^{-1}).
 \label{eq:laurent-square}
\end{equation}
If $m$ and $n$ are the smallest and largest exponents with nonzero coefficients in
$f$, then the smallest and largest exponents in $f^2$ are $2m$ and $2n$; their
coefficients cannot cancel because each is the square of an extreme coefficient. The
right-hand side of Eq.~\eqref{eq:laurent-square} has extreme exponents $-1$ and $1$
when $t\ne0$, which is impossible because $2m$ and $2n$ are even. Thus $d$ is not a
finite Fourier polynomial. Equation~\eqref{eq:projected-range} has diagonal entries
containing $-d(k)/2$ plus finite Fourier terms, so the embedded projected operator is not
finite range.
\end{proof}

\subsection{Flat-parent residual--locality identity}
\label{app:proof-flat-parent}

\begin{proof}[Proof of Theorem~\ref{thm:flat-parent}]
A constant real rotation diagonalizes $D$ without changing any Frobenius norm, Fourier
weight, or minimization in the statement, so take $D=\delta\sigma_z$ with
$\delta\ne0$. Conjugation through a real rotation gives
\begin{equation}
 B_a^{(d)}(\mathbf q)
 =\delta\left[
 \cos\!\left(2\theta_a^{(d)}(\mathbf q)\right)\sigma_z
 -\sin\!\left(2\theta_a^{(d)}(\mathbf q)\right)\sigma_x
 \right].
 \label{eq:flat-rotated-pauli}
\end{equation}
Using the normalized integral representation
\begin{equation}
 J_0(x)=\int_{-1/2}^{1/2}e^{ix\sin(2\pi q)}\,dq
 \label{eq:bessel-integral}
\end{equation}
and the product structure of $\T^d$,
\begin{align}
 \int_{\T^d}e^{2i\theta_a^{(d)}(\mathbf q)}\,d\mathbf q
 &=\prod_{\alpha=1}^d
   \int_{-1/2}^{1/2}
   e^{i(2a/d)\sin(2\pi q_\alpha)}\,dq_\alpha\\
 &=J_0\!\left(\frac{2a}{d}\right)^d=c_d(a).
\end{align}
Global inversion $\mathbf q\mapsto-\mathbf q$ makes the sine average in
Eq.~\eqref{eq:flat-rotated-pauli} vanish, proving
Eq.~\eqref{eq:onsite-average}. Moreover,
$0\le2a/d\le\pi/4$, and the real part of the integrand in
Eq.~\eqref{eq:bessel-integral} is at least $\cos(2a/d)>0$; hence $c_d(a)>0$.

For any constant matrix $C$, the zero-mean property of
$B_a^{(d)}-B_{\mathbf0}^{(d)}$ gives the orthogonal decomposition
\begin{equation}
 \int_{\T^d}\norm{C-B_a^{(d)}(\mathbf q)}_F^2\,d\mathbf q
 =\int_{\T^d}\norm{B_a^{(d)}(\mathbf q)-B_{\mathbf0}^{(d)}}_F^2\,d\mathbf q
  +\norm{C-B_{\mathbf0}^{(d)}}_F^2.
 \label{eq:constant-projection}
\end{equation}
The constant $B_{\mathbf0}^{(d)}=c_d(a)D$ belongs to the allowed family in
Eq.~\eqref{eq:flat-operator-loss}: choose $s=0$, $\lambda=c_d(a)$, and $\phi=0$.
It is therefore the best constant restricted alignment. Equivalently, the normalized
maximum correlation with one global rotation is $c_d(a)$.

Parseval's identity applied entrywise yields
\begin{equation}
 \int_{\T^d}\norm{B_a^{(d)}-B_{\mathbf0}^{(d)}}_F^2
 =\sum_{\mathbf R\ne\mathbf0}\norm{B_{\mathbf R}^{(d)}}_F^2.
\end{equation}
Pointwise unitary invariance gives
$\norm{B_a^{(d)}(\mathbf q)}_F=\norm D_F$. A second use of Parseval, together with
Eq.~\eqref{eq:onsite-average}, therefore gives
\begin{equation}
 \sum_{\mathbf R\ne\mathbf0}\norm{B_{\mathbf R}^{(d)}}_F^2
 =\norm D_F^2-\norm{B_{\mathbf0}^{(d)}}_F^2
 =\norm D_F^2\left[1-c_d(a)^2\right].
\end{equation}
Dividing by $2E_{\mathrm{ref}}^2$ proves Eqs.~\eqref{eq:flat-residual} and
\eqref{eq:flat-tail}. Finally, fixing $s=0$ and $\lambda=1$ leaves the non-onsite
weight plus the onsite mismatch:
\begin{align}
 \inf_\phi L_{\mathrm{op},\phi}^2(0,1)
 &=\kappa[1-c_d(a)^2]+\kappa[1-c_d(a)]^2\\
 &=2\kappa[1-c_d(a)],
\end{align}
which proves Eq.~\eqref{eq:flat-exact-spectrum}.
\end{proof}

\subsection{Orthogonal shell coefficients and factorial tail}
\label{app:proof-orthogonal-tail}

\begin{proof}[Proof of Theorem~\ref{thm:orthogonal-fourier-tail}]
The Jacobi--Anger identity gives
\begin{equation}
 e^{2i\theta_a^{(d)}(\mathbf q)}
 =\prod_{\alpha=1}^d
  \sum_{n\in\mathbb Z}J_n(z)e^{2\pi inq_\alpha}
 =\sum_{\mathbf R\in\mathbb Z^d}
  j_{\mathbf R}(z)e^{2\pi i\mathbf q\cdot\mathbf R}.
\end{equation}
Because $J_{-n}(z)=(-1)^nJ_n(z)$, the cosine coefficient vanishes for odd
$p(\mathbf R)$ and equals $j_{\mathbf R}(z)$ for even $p(\mathbf R)$; the sine
coefficient vanishes for even parity and equals $-ij_{\mathbf R}(z)$ for odd parity.
Substitution into Eq.~\eqref{eq:flat-rotated-pauli} proves
Eq.~\eqref{eq:orthogonal-fourier-block}. Since both Pauli matrices have squared
Frobenius norm two,
\begin{equation}
 \norm{B_{\mathbf R}^{(d)}}_F^2
 =\norm D_F^2\prod_{\alpha=1}^dJ_{R_\alpha}(z)^2.
 \label{eq:block-power-product}
\end{equation}
Equation~\eqref{eq:shell-power} follows by grouping a complete shell. Parseval applied
to the unit-modulus function $e^{iz\sin(2\pi q)}$ gives
$\sum_{n\in\mathbb Z}J_n(z)^2=1$; taking the $d$-fold product proves
Eq.~\eqref{eq:shell-partition}.

If $\norm{\mathbf R}_2^2>S$, at least one coordinate obeys $|R_\alpha|\ge L$.
The union bound and the one-dimensional Bessel partition therefore give
\begin{equation}
 \sum_{\norm{\mathbf R}_2^2>S}
 \prod_{\alpha=1}^dJ_{R_\alpha}(z)^2
 \le 2d\sum_{n=L}^{\infty}|J_n(z)|^2.
 \label{eq:tail-union}
\end{equation}
The defining Bessel series and $(n+m)!\ge n!m!$ imply, for $n\ge0$,
\begin{equation}
 |J_n(z)|
 \le \frac{x^n}{n!}\sum_{m=0}^{\infty}\frac{x^{2m}}{m!}
 =\frac{x^n}{n!}e^{x^2}.
\end{equation}
Moreover, $(L+m)!\ge L!(L+1)^m$, so when $x<L+1$,
\begin{equation}
 \sum_{n=L}^{\infty}\frac{x^{2n}}{(n!)^2}
 \le \frac{x^{2L}}{(L!)^2}
 \sum_{m=0}^{\infty}\left(\frac{x^2}{(L+1)^2}\right)^m
 =\frac{x^{2L}}
 {(L!)^2\left[1-x^2/(L+1)^2\right]}.
\end{equation}
Combining this estimate with Eq.~\eqref{eq:tail-union}, using
$e^{2x^2}=e^{z^2/2}$ and Eq.~\eqref{eq:block-power-product}, proves
Eq.~\eqref{eq:orthogonal-factorial-tail}.
\end{proof}

\subsection{Monotonicity and small-amplitude dimension law}
\label{app:proof-dimension-law}

\begin{proof}[Proof of Corollary~\ref{cor:dimension-law}]
For $x>0$, $J_0'(x)=-J_1(x)$. On $0<x\le\pi/4$, the alternating series
\begin{equation}
 J_1(x)=\sum_{m=0}^{\infty}
 \frac{(-1)^m}{m!(m+1)!}\left(\frac{x}{2}\right)^{2m+1}
\end{equation}
has strictly decreasing term magnitudes because their successive ratio is at most
$\pi^2/128<1$. Hence $J_1(x)>0$. Together with the positivity already proved for
$J_0(2a/d)$, this shows
\begin{equation}
 c_d'(a)=-2J_0(2a/d)^{d-1}J_1(2a/d)<0
\end{equation}
for $a>0$. The threshold monotonicity follows immediately. Finally,
$J_0(x)=1-x^2/4+x^4/64+O(x^6)$. Substituting $x=2a/d$, raising to the $d$th power, and
expanding the two threshold formulas proves
Eqs.~\eqref{eq:correlation-small-a}--\eqref{eq:exact-threshold-small-a}.
\end{proof}

\subsection{Bond-orbit flat-parent comparator}
\label{app:proof-bond-orbit}

\begin{proof}[Proof of Theorem~\ref{thm:bond-orbit}]
Equation~\eqref{eq:bond-orbit-phase} is odd under
$\mathbf q\mapsto-\mathbf q$, so the imaginary part of
Eq.~\eqref{eq:bond-orbit-correlation} integrates to zero. Since
$|2\theta_{a,\mathcal B}|\le\pi/2$, its cosine is nonnegative; it is positive in a
neighborhood of $\mathbf q=\mathbf0$, where the phase vanishes. Hence the integral is
strictly positive. Diagonalizing $D$ by one constant real rotation and repeating
Eq.~\eqref{eq:flat-rotated-pauli} gives
$D_{\mathbf0}^{\mathcal B}=c_{\mathcal B,\boldsymbol\eta}D$. The constant-projection
identity in Eq.~\eqref{eq:constant-projection} and entrywise Parseval then prove the two
residual formulas exactly as in Theorem~\ref{thm:flat-parent}.

For the explicit comparator, apply the Jacobi--Anger expansion to every bond harmonic:
\begin{equation}
 e^{i(2a\eta_j/W)\sin(2\pi\mathbf q\cdot\mathbf b_j)}
 =\sum_{n_j\in\Z}
 J_{n_j}\!\left(\frac{2a\eta_j}{W}\right)
 e^{2\pi i n_j\mathbf q\cdot\mathbf b_j}.
\end{equation}
The finite product of these absolutely convergent series may be integrated termwise.
Torus orthogonality retains exactly the tuples satisfying
$\sum_jn_j\mathbf b_j=\mathbf0$, which proves
Eq.~\eqref{eq:bond-orbit-bessel-sum}. If the integer kernel of the bond matrix is
trivial, only $\mathbf n=\mathbf0$ remains, giving
Eq.~\eqref{eq:independent-bond-product}.
\end{proof}

\subsection{Constrained-series complete-shell certificate}
\label{app:proof-constrained-tail}

\begin{proof}[Proof of Theorem~\ref{thm:constrained-tail}]
Jacobi--Anger expansion and torus orthogonality give
Eq.~\eqref{eq:constrained-coefficient}. The symmetry
$e^{2i\theta(-\mathbf q)}=\overline{e^{2i\theta(\mathbf q)}}$ makes every
$g_{\mathbf R}$ real. Splitting the cosine and sine coefficients into
$(g_{\mathbf R}\pm g_{-\mathbf R})/2$ proves
Eq.~\eqref{eq:constrained-matrix-block}. Hence
\begin{equation}
 \frac{\norm{D_{\mathbf R}^{\mathcal B}}_F^2}{\norm D_F^2}
 =\frac{g_{\mathbf R}^2+g_{-\mathbf R}^2}{2}.
\end{equation}
Parseval for the unit-modulus scalar phase gives
$\sum_{\mathbf R}g_{\mathbf R}^2=1$. Because $\mathcal S_S$ is inversion complete,
the exact matrix tail normalized by $\norm D_F^2$ is
$1-\sum_{\mathbf R\in\mathcal S_S}g_{\mathbf R}^2$. The reverse triangle inequality
$|g_{\mathbf R}|\ge\max\{0,|\widetilde g_{\mathbf R}|-\epsilon_{\mathbf R}\}$
proves Eq.~\eqref{eq:constrained-tail-certificate}.

Solving the triangular constraint gives
$(n_1,n_2,n_3)=(r_1-\ell,r_2+\ell,\ell)$, which proves
Eq.~\eqref{eq:triangular-coefficient}. Taking
$(u,v,w)=(n_4,n_5,n_6)$ in the FCC constraint gives
\begin{equation}
 (n_1,n_2,n_3)=(r_1+u+v,r_2-u+w,r_3-v-w),
\end{equation}
which proves Eq.~\eqref{eq:fcc-coefficient}. Finally, the defining Bessel series gives
$|J_n(z)|\le e^{x^2}x^{|n|}/|n|!$. Summing it over all integers gives $A_0$; summing
over $|n|>K$ and bounding the remaining exponential-series tail geometrically gives
$A_K$. In the triangular remainder the other two Bessel factors have magnitude at most
one. In the FCC remainder, a union bound over the three free indices gives
$3A_KA_0^2$.
\end{proof}
